\section{Research Experience}

\entry{Graduate Researcher}{2022 -- present}
  {Parallel Programming Laboratory, University of Illinois Urbana-Champaign}{Urbana, IL}
\begin{itemize}
  \item Primary designer of \href{https://github.com/charmplusplus/reconverse}{Reconverse},
    a scheduling and communication layer replacing Charm++'s 30-year-old Converse, with
    lottery-queue scheduling, queue registration, and an LCI machine layer; up to
    2$\times$ speedups on large-scale applications
  \item Developed ACIC, an adaptive asynchronous shortest-path algorithm for high-diameter
    graphs with tens of billions of vertices, up to two orders of magnitude faster than
    existing distributed graph frameworks
  \item Accelerated friends-of-friends cluster finding for cosmological simulations by
    combining tree-traversal and union-find libraries with fine-grained optimizations
    for highly imbalanced datasets
  \item Implemented the Ray Core API on top of
    \href{https://github.com/charmplusplus/charm4py}{Charm4Py}, letting unmodified Ray
    programs use dynamic load balancing by mapping Ray actors to Charm4Py chares;
    3--4$\times$ speedups on Ray benchmarks over the standard Ray runtime
  \cvonly{\item Built a library of graph algorithms for CharmTyles, an interactive
    computing interface to Charm++ aimed at data scientists}
\end{itemize}
\cvonly{%
\entrygap

\entry{Undergraduate Researcher, Julia on AMD GPUs}{June 2021 -- June 2022}
  {NUCAR Lab, Northeastern University}{Boston, MA}
\begin{itemize}
  \item Designed an interface that translates Julia GPU code to HIP for AMD's ROCm
    platform, allowing Oceananigans, an ocean climate model written in Julia, to run on
    AMD GPUs
\end{itemize}
\entrygap

\entry{Pi-Tiles}{Jan.\ 2019 -- Apr.\ 2019}{NUCAR Lab, Northeastern University}{Boston, MA}
\begin{itemize}
  \item Implemented blocked matrix multiplication in MPI across a cluster of Raspberry
    Pis, 10$\times$ faster than the naive method
\end{itemize}}
